\newcommand{\figuraAngulosContacto}{%
\begin{figure}[H]
\centering
\begin{tikzpicture}[>=Latex,font=\small,line cap=round,line join=round]
  \node[font=\bfseries] at (3,4.8) {(a) Líquido mojante};
  \draw[very thick] (0.6,0.7)--(5.4,0.7) (1.1,0.7)--(1.1,4.2) (4.9,0.7)--(4.9,4.2);
  \fill[blue!15] (1.18,0.78)--(4.82,0.78)--(4.82,2.0)
    .. controls (4.0,1.55) and (2.0,1.55) .. (1.18,2.0)--cycle;
  \draw[very thick,blue!70!black] (1.18,2.0)
    .. controls (2.0,1.55) and (4.0,1.55) .. (4.82,2.0);
  \draw[->,red!75!black] (1.18,2.0) arc[start angle=90,end angle=25,radius=0.65];
  \node[red!75!black] at (1.75,2.45) {$\theta<90^\circ$};

  \node[font=\bfseries] at (9,4.8) {(b) Líquido no mojante};
  \draw[very thick] (6.6,0.7)--(11.4,0.7) (7.1,0.7)--(7.1,4.2) (10.9,0.7)--(10.9,4.2);
  \fill[gray!18] (7.18,0.78)--(10.82,0.78)--(10.82,1.65)
    .. controls (10.0,2.15) and (8.0,2.15) .. (7.18,1.65)--cycle;
  \draw[very thick,black!65] (7.18,1.65)
    .. controls (8.0,2.15) and (10.0,2.15) .. (10.82,1.65);
  \draw[->,red!75!black] (7.18,1.65) arc[start angle=90,end angle=-35,radius=0.65];
  \node[red!75!black] at (7.85,2.45) {$\theta>90^\circ$};
\end{tikzpicture}
\caption{Efecto del ángulo de contacto. Un líquido mojante asciende junto a
la pared y forma un menisco cóncavo; uno no mojante desciende y forma un
menisco convexo.}
\label{fig:angulos-contacto}
\end{figure}%
}